% ============================================================
%  Manual de Instalación — Churn Dashboard (Olist)
% ============================================================
\documentclass[11pt]{article}

\usepackage[margin=1in]{geometry}
\usepackage[T1]{fontenc}
\usepackage[utf8]{inputenc}
\usepackage{lmodern}
\usepackage{microtype}
\usepackage{graphicx}
\usepackage{float}
\usepackage{listings}
\usepackage{xcolor}
\usepackage{array}
\usepackage{booktabs}
\usepackage{enumitem}
\usepackage{hyperref}
\hypersetup{colorlinks=true,linkcolor=blue,urlcolor=blue}
\usepackage{bookmark}
\usepackage{url}

% Estilo para bloques de código bash
\lstdefinestyle{bash}{
  language=bash,
  basicstyle=\ttfamily\small,
  backgroundcolor=\color{gray!10},
  frame=single,
  rulecolor=\color{gray!40},
  breaklines=true,
  breakatwhitespace=false,
  showstringspaces=false,
  keywordstyle=\color{blue},
  commentstyle=\color{gray},
  stringstyle=\color{teal},
  numbers=none,
}

\lstset{style=bash}

\title{Manual de Instalación\\
\large Churn Dashboard --- Predicción de Abandono de Clientes Olist}
\author{Equipo Microproyecto}
\date{Marzo 2026}

\begin{document}
\maketitle
\tableofcontents
\newpage

% ============================================================
\section{Descripción general}

Este manual describe los pasos para instalar y poner en marcha el
\textbf{Churn Dashboard} en una instancia \textbf{Ubuntu 22.04} (recomendado
en AWS EC2). Al finalizar, el tablero estará accesible en el puerto
\texttt{3000} del servidor y administrado por PM2.

\subsection*{Componentes instalados}

\begin{itemize}[itemsep=2pt]
  \item \textbf{Frontend}: aplicación React/Vite compilada y servida por
        Express desde el directorio \texttt{dist/}.
  \item \textbf{Backend}: servidor Express (Node.js) en
        \texttt{server/index.js}, endpoint \texttt{POST /api/predict}.
  \item \textbf{Inferencia ML}: script Python \path{server/predict_joblib.py}
        con el modelo \path{best_model_logreg.joblib}.
  \item \textbf{Gestor de procesos}: PM2 mantiene el servicio activo tras
        reinicios.
\end{itemize}

\subsection*{Repositorio}

\begin{center}
  \url{https://github.com/Jason-Morrill/Microproyecto_proyectogrado}
\end{center}

% ============================================================
\section{Requisitos previos}

\begin{center}
\begin{tabular}{lll}
\toprule
\textbf{Componente} & \textbf{Versión mínima} & \textbf{Notas} \\
\midrule
Ubuntu / Debian  & 22.04 LTS & Recomendado EC2 \\
Node.js          & 18.x      & LTS recomendado \\
npm              & 9.x       & Se instala con Node.js \\
Python           & 3.12      & Con \texttt{venv} y \texttt{pip} \\
PM2              & 5.x       & Instalado globalmente vía npm \\
Git              & 2.x       & Para clonar el repositorio \\
\bottomrule
\end{tabular}
\end{center}

\noindent\textbf{Puertos:} el Security Group (EC2) o firewall debe permitir
tráfico entrante en el puerto \texttt{3000} (TCP) desde las IPs que accederán
al tablero.

% ============================================================
\section{Instalación paso a paso}

\subsection{Paso 1 — Actualizar el sistema}

\begin{lstlisting}
sudo apt update && sudo apt upgrade -y
\end{lstlisting}

\subsection{Paso 2 — Instalar Node.js y npm}

\begin{lstlisting}
sudo apt install -y nodejs npm
node -v    # debe mostrar v18.x o superior
npm -v
\end{lstlisting}

Si la versión de Node.js disponible en \texttt{apt} es antigua, instale la
versión LTS mediante NodeSource:

\begin{lstlisting}
curl -fsSL https://deb.nodesource.com/setup_18.x | sudo -E bash -
sudo apt install -y nodejs
node -v
\end{lstlisting}

\subsection{Paso 3 — Instalar PM2 de forma global}

\begin{lstlisting}
sudo npm install -g pm2
pm2 -v
\end{lstlisting}

\subsection{Paso 4 — Configurar arranque automático de PM2}

\begin{lstlisting}
pm2 startup systemd -u ubuntu --hp /home/ubuntu
\end{lstlisting}

El comando anterior imprimirá otro comando con \texttt{sudo}. Ejecútelo
exactamente como aparezca en pantalla. A continuación:

\begin{lstlisting}
pm2 save
\end{lstlisting}

\subsection{Paso 5 — Instalar Python 3.12 + venv + pip}

\begin{lstlisting}
sudo apt install -y python3.12-venv python3-pip
python3 --version     # debe mostrar 3.12.x
python3 -m pip --version
\end{lstlisting}

\subsection{Paso 6 — Clonar el repositorio}

\begin{lstlisting}
cd /home/ubuntu
git clone https://github.com/Jason-Morrill/Microproyecto_proyectogrado \
    Microproyecto_proyectogrado
cd /home/ubuntu/Microproyecto_proyectogrado
\end{lstlisting}

Si el repositorio ya existe, actualícelo:

\begin{lstlisting}
cd /home/ubuntu/Microproyecto_proyectogrado
git pull --ff-only
\end{lstlisting}

\subsection{Paso 7 — Ejecutar el script de despliegue automático}

\begin{lstlisting}
cd /home/ubuntu/Microproyecto_proyectogrado/Tablero_final
bash deploy_churn_dashboard.sh
\end{lstlisting}

El script realiza automáticamente:

\begin{enumerate}[itemsep=2pt]
  \item Instalación de dependencias del frontend (\texttt{npm install}).
  \item \texttt{npm run build} para compilar la aplicación Vite en \texttt{dist/}.
  \item Instalación de dependencias del backend (\texttt{npm install} en
        \texttt{server/}).
  \item Creación (o reparación) del entorno virtual Python en
        \texttt{server/.venv}.
  \item Instalación de dependencias Python desde
        \texttt{server/requirements.txt}.
  \item Inicio (o reinicio) del proceso \texttt{churn\_dashboard} en PM2.
  \item \texttt{pm2 save} para persistir el proceso.
\end{enumerate}

Al finalizar correctamente verá en consola:

\begin{lstlisting}
[PM2] Process churn_dashboard online
\end{lstlisting}

% ============================================================
\section{Verificación de la instalación}

\subsection{Estado del proceso PM2}

\begin{lstlisting}
pm2 status
\end{lstlisting}

El proceso \texttt{churn\_dashboard} debe aparecer en estado \textbf{online}.

\subsection{Prueba del endpoint de predicción}

\begin{lstlisting}
curl -sS -X POST http://127.0.0.1:3000/api/predict \
  -H "Content-Type: application/json" \
  -d '{
    "recency": 32,
    "frequency": 5,
    "monetary": 500,
    "avg_review_score": 4,
    "avg_delivery_days": 10,
    "avg_late_days": 0,
    "avg_num_items": 2,
    "avg_price_sum": 100,
    "avg_freight_sum": 15,
    "customer_state": "SP"
  }'
\end{lstlisting}

Respuesta esperada:

\begin{lstlisting}
{"churnProbability": 0.1234}
\end{lstlisting}

\subsection{Acceso desde el navegador}

Abra un navegador y navegue a:

\begin{center}
  \texttt{http://<IP-PUBLICA-EC2>:3000}
\end{center}

Debe cargarse la interfaz del tablero con el formulario lateral de entrada
de datos del cliente.

% ============================================================
\section{Comandos de operación}

\begin{center}
\begin{tabular}{>{}p{4.2cm}p{9.5cm}}
\toprule
\textbf{Acción} & \textbf{Comando} \\
\midrule
Ver estado del proceso &
  \texttt{pm2 status} \\[4pt]
Ver logs en tiempo real &
  \texttt{pm2 logs churn\_dashboard -{}-lines 100} \\[4pt]
Ver logs de aplicación &
  \small\path{tail -f /home/ubuntu/Microproyecto_proyectogrado/Tablero_final/server/logs/server.log} \\[4pt]
Reiniciar tras cambios &
  \texttt{pm2 restart churn\_dashboard -{}-update-env} \\[4pt]
Eliminar y recrear proceso &
  \texttt{pm2 delete churn\_dashboard} \\
  &
  \texttt{pm2 start .../server/index.js -{}-name churn\_dashboard} \\
  &
  \texttt{pm2 save} \\
\bottomrule
\end{tabular}
\end{center}

% ============================================================
\section{Resolución de problemas frecuentes}

\subsection{Error \texttt{ERR\_HTTP\_HEADERS\_SENT}}

\textbf{Causa:} doble respuesta HTTP generada en el servidor.

\textbf{Solución:}
\begin{lstlisting}
cd /home/ubuntu/Microproyecto_proyectogrado
git pull --ff-only
pm2 restart churn_dashboard --update-env
pm2 logs churn_dashboard --lines 100
\end{lstlisting}

\subsection{Error: \texttt{pandas}, \texttt{joblib} o \texttt{sklearn} no encontrado}

\textbf{Causa:} el entorno virtual Python no tiene las dependencias instaladas
o el proceso PM2 no usa el Python de \texttt{.venv}.

\textbf{Verificación:}
\begin{lstlisting}
/home/ubuntu/Microproyecto_proyectogrado/Tablero_final/server/.venv/bin/python \
  -c "import pandas, joblib, sklearn; print('ok')"
\end{lstlisting}

\textbf{Solución si falla:}
\begin{lstlisting}
cd /home/ubuntu/Microproyecto_proyectogrado/Tablero_final/server
source .venv/bin/activate
python -m pip install -r requirements.txt
deactivate
pm2 restart churn_dashboard --update-env
\end{lstlisting}

\subsection{No se registran requests en \texttt{server.log}}

\begin{enumerate}[itemsep=2pt]
  \item Verifique la ruta exacta del log con \texttt{pm2 show churn\_dashboard}.
  \item Genere tráfico con \texttt{curl} (ver §4.2).
  \item Confirme que el proceso usa el script correcto:
        \texttt{pm2 show churn\_dashboard | grep script}.
\end{enumerate}

\subsection{La página no carga desde el navegador}

\begin{enumerate}[itemsep=2pt]
  \item Confirme que \texttt{pm2 status} muestra \texttt{churn\_dashboard}
        en estado \texttt{online}.
  \item Verifique que el puerto 3000 esté abierto en el Security Group de EC2
        (Inbound rules: TCP 3000 desde su IP o \texttt{0.0.0.0/0}).
  \item Pruebe localmente desde el propio servidor con
        \texttt{curl http://127.0.0.1:3000}.
\end{enumerate}

% ============================================================
\section{Seguridad y buenas prácticas}

\begin{itemize}[itemsep=4pt]
  \item \textbf{Security Group EC2}: restrinja el puerto 3000 solo a las IPs
        que necesiten acceso. Para producción, coloque un balanceador de carga
        (ALB) con HTTPS y no exponga el puerto 3000 directamente.
  \item \textbf{Secretos}: no suba archivos \texttt{.env} con credenciales al
        repositorio. Use AWS Secrets Manager o variables de entorno del sistema.
  \item \textbf{Modelo}: versione el artefacto \texttt{best\_model\_logreg.joblib}
        en MLflow y no lo sobreescriba sin incrementar la versión.
  \item \textbf{Instancia EC2}: mantenga la instancia \textbf{detenida}
        (Stopped) cuando no esté en uso. \textbf{No la termine (Terminate)},
        ya que se perdería la configuración instalada.
\end{itemize}

\end{document}
